\documentclass[11pt]{article}
\usepackage[utf8]{inputenc}
\usepackage[T1]{fontenc}
\usepackage{lmodern}
\usepackage{amsmath,amssymb}
\usepackage{hyperref}

\title{Rinnovo Binary (RNB): Artifact Format and Runtime}
\author{}
\date{}

\begin{document}
\maketitle

\section{Overview}

Rinnovo Binary (RNB) is a single-file container format and runtime
interface for biological datasets. An \emph{RNB artifact} is a compiled,
portable representation of a biological dataset and its semantics:
once produced, the artifact can be opened by a runtime to answer
queries without requiring the original workflow, intermediate files,
or ad-hoc glue code.

Conceptually, RNB sits between file formats and databases:
\begin{itemize}
  \item like a file format, an RNB artifact is self-contained, easily
        distributed with publications and collaborations, and can be
        version-controlled;
  \item like a database, it offers a typed, segment-based layout
        designed for memory-mapped access and deterministic query
        performance.
\end{itemize}

The goals of the format and runtime are:
\begin{itemize}
  \item \textbf{Semantic stability}: the logical model (objects,
        relationships, and query kernels) should be stable even as
        physical layouts evolve;
  \item \textbf{Determinism}: compiling the same semantic input must
        produce bit-for-bit identical artifacts;
  \item \textbf{Performance}: common analyses should be expressible as
        simple, bounded-time operations over memory-mapped segments;
  \item \textbf{Portability}: the on-disk format should be simple
        enough to reimplement in multiple languages with a small,
        auditable code footprint.
\end{itemize}

\section{File Layout and Manifest}

This section describes the top-level binary layout of an RNB artifact
and the role of the manifest segment as the semantic contract between
artifacts and runtimes.

\subsection{File Layout}

Every RNB artifact has the following high-level structure:
\begin{verbatim}
[RNB HEADER][SEGMENT DATA ...][SEGMENT DIRECTORY]
\end{verbatim}

\paragraph{Header.}
The header is a fixed-size structure containing:
\begin{itemize}
  \item magic bytes: \verb|RNB\0|;
  \item version: 16-bit major and minor fields;
  \item directory offset: 64-bit little-endian file offset to the
        segment directory;
  \item directory length: 64-bit little-endian length in bytes of the
        segment directory.
\end{itemize}

The header is validated on open (magic and version), after which the
runtime locates and parses the directory.

\paragraph{Segment directory.}
The segment directory is an array of fixed-width entries. Each entry
describes a single segment:
\begin{itemize}
  \item \verb|segment_id| (u32): a caller-assigned identifier for the
        segment instance;
  \item \verb|segment_type| (u32): the logical type of the segment,
        drawn from an enumerated \texttt{SegmentType};
  \item \verb|offset| (u64): byte offset of the segment payload
        relative to the beginning of the file;
  \item \verb|length| (u64): length in bytes of the segment payload;
  \item \verb|checksum| (u64): a 64-bit FNV-1a checksum over the
        segment payload.
\end{itemize}

Unknown \verb|segment_type| values are required to be skippable:
readers can safely ignore segments they do not understand by using
only the directory metadata.

\subsection{Manifest Segment}

The manifest segment (segment type \texttt{Manifest}) is required and
must appear exactly once in every RNB artifact. It is the primary
semantic descriptor for the file: it declares which segment types are
required, which query kernels are supported, and global configuration
parameters such as maximum chunk size.

\paragraph{Binary layout.}
The manifest segment is a versioned, little-endian structure:
\begin{itemize}
  \item magic: \verb|MNF\0| (4 bytes);
  \item version: 16-bit major and minor fields (currently 0 and 1);
  \item \verb|flags| (u32): reserved for future format-level options;
  \item \verb|required_segment_count| (u32);
  \item \verb|supported_kernel_count| (u32);
  \item \verb|max_chunk_bytes| (u32): an advisory upper bound on chunk
        sizes for streaming or incremental processing;
  \item \verb|reserved| (u32): must be zero;
  \item \verb|required_segments|: an array of
        \verb|required_segment_count| 32-bit values, each encoding a
        \texttt{SegmentType};
  \item \verb|supported_kernels|: an array of
        \verb|supported_kernel_count| 32-bit values, each encoding a
        \texttt{QueryKernel}.
\end{itemize}

Unknown \texttt{SegmentType} or \texttt{QueryKernel} identifiers are
treated as hard errors when parsing the manifest: they indicate that
the artifact requires capabilities beyond what the current runtime
understands.

\paragraph{Required segments.}
The \verb|required_segments| list expresses a semantic requirement:
for each listed \texttt{SegmentType}, at least one directory entry
with a matching \verb|segment_type| value MUST be present. The RNB
container loader enforces this invariant when opening an artifact: if
any required segment type declared in the manifest is absent from the
directory, the artifact is rejected as invalid.

This mechanism decouples the manifest from specific segment IDs while
still allowing the runtime to rely on the presence of certain segment
types (for example, a global string dictionary or an object table) in
order to implement higher-level query kernels.

\paragraph{Supported kernels.}
The \verb|supported_kernels| list advertises the set of query kernels
that the artifact is capable of serving. Each kernel is identified by
an integer \texttt{QueryKernel} value (e.g.,
\texttt{GetObjectById = 1}). The mapping from kernel identifiers to
behavior is defined at the runtime level, but the manifest ensures
that a client can discover, at open time, which kernels are supported
for a given artifact.

\section{Segments for Strings and Objects}

This section summarizes the core segments used to describe strings and
objects in RNB v0.1. Each segment has a small, fixed-width header and
is designed to support constant-time offset computations under
memory-mapping.

\subsection{String Dictionary Segment}

The String Dictionary segment (\texttt{SegmentType::StringDict = 2})
stores a compact dictionary of UTF-8 strings used elsewhere in the
artifact. Its layout is:
\begin{itemize}
  \item magic: \verb|SDCT| (4 bytes);
  \item version: 16-bit major and minor fields (currently 0 and 1);
  \item \verb|string_count| (u32);
  \item \verb|blob_len| (u32): total number of bytes in the UTF-8 blob;
  \item \verb|offsets|: an array of \verb|string_count + 1| 32-bit
        offsets into the blob;
  \item \verb|blob|: \verb|blob_len| bytes of concatenated UTF-8
        string data.
\end{itemize}

The \verb|offsets[i]| and \verb|offsets[i+1]| entries delimit the byte
range of string $i$ within the blob. Offsets MUST start at 0, MUST be
non-decreasing, and the final offset MUST equal \verb|blob_len|.
Strings MUST be valid UTF-8. This layout supports constant-time
lookup of a string by its integer identifier.

\subsection{Object Table Segment}
\label{sec:object-table}

The Object Table segment (\texttt{SegmentType::ObjectTable = 3})
provides a minimal, fixed-width mapping from object identifiers to
basic metadata. It is laid out as:
\begin{itemize}
  \item magic: \verb|OBT\0| (4 bytes);
  \item version: 16-bit major and minor fields (currently 0 and 1);
  \item \verb|object_count| (u32);
  \item \verb|reserved| (u32): must be zero;
  \item rows: \verb|object_count| repetitions of:
  \begin{itemize}
    \item \verb|type_sid| (u32): StringDict ID for the object's
          type/kind;
    \item \verb|name_sid| (u32): StringDict ID for a primary
          label/name;
    \item \verb|flags| (u32): reserved for future use.
  \end{itemize}
\end{itemize}

By construction, the $i$-th row (0-based) corresponds to
\verb|object_id = i|. This allows constant-time access to object
metadata by identifier using simple offset arithmetic.

\section{Logical Object Model}

The logical object model exposed by the runtime is a thin semantic
layer over the Object Table and String Dictionary segments.

Each object is identified by an unsigned 32-bit \verb|object_id|,
which is defined to be the row index in the Object Table segment. The
runtime presents a stable view:
\begin{verbatim}
struct Object {
    id: u32;
    type_sid: u32;
    name_sid: u32;
    flags: u32;
}
\end{verbatim}

where \verb|type_sid| and \verb|name_sid| are StringDict identifiers
into the global dictionary. Additional segments such as the attribute
table, relation table, or numeric matrices may attach further
information to objects by referring to their \verb|object_id|, but the
core identity and basic metadata of an object are derived solely from
the Object Table.

This separation between on-disk segments and the logical object view
allows the physical layout to evolve (for example, using columnar
storage or indices) without changing the semantics seen by callers of
the runtime.

\end{document}
